
\documentclass[12pt]{report}

\usepackage{graphicx}
\usepackage{amsmath}
\usepackage{cite}

\Large\title{\textbf{The Solar System}}
\author{Anulekha}
\date{May 28, 2026}


\begin{document}


\maketitle
\tableofcontents

%---------------- REPORT ----------------%

\chapter{Introduction}

\section{What is the Solar System}

The Solar System is the Sun and all the celestial objects that revolve around it due to the force of gravity. It includes eight planets, dwarf planets, moons, asteroids, comets, and other small celestial bodies.

\section{Importance of Studying the Solar System}

Studying the Solar System helps scientists understand the origin of planets, the evolution of Earth, and the possibility of life beyond our planet. Space exploration also leads to technological advancements.

%--------------------------------------------%

\chapter{Methodology}

\section{Components of the Solar System}

The Solar System consists of the Sun, eight planets, dwarf planets, moons, asteroids, comets, and interplanetary dust.

\section{Classification of Planets}

The planets are classified into two groups:

\begin{itemize}
\item Inner Planets (Mercury, Venus, Earth, Mars)
\item Outer Planets (Jupiter, Saturn, Uranus, Neptune)
\end{itemize}

%--------------------------------------------%

\chapter{Results}

\section{Overview of the Eight Planets}

The Solar System officially contains eight planets that orbit the Sun. Figure~\ref{fig:solar} shows the Solar System.

\begin{figure}[htbp]
\centering
\vspace{1cm}
\includegraphics[width=0.2\textwidth]{solar.jpeg}
\caption{Solar System}
\label{fig:solar}
\end{figure}

%--------------------------------------------%

\chapter{Conclusion}

\section{Summary}

The motion of planets around the Sun is governed by Newton's Law of Universal Gravitation. Equation~\ref{eq:formula} represents this relationship.

\begin{equation}
F = G\frac{m_1m_2}{r^2}
\label{eq:formula}
\end{equation}

%--------------------------------------------%

\chapter{References}

\begin{itemize}
\item NASA Solar System Exploration Website
\item Encyclopaedia Britannica -- Solar System
\item European Space Agency (ESA) -- Solar System Information
\end{itemize}

%================ ARTICLE ===================%

\newpage

\section{Article Section}

These results demonstrate that gravity plays a significant role in maintaining the stability of the Solar System. Every planet follows a predictable orbit because of gravitational attraction.

\begin{figure}[htbp]
\centering
\includegraphics[width=0.5\textwidth]{space.jpeg}
\caption{Spacecraft Trajectory in the Solar System}
\label{fig:space}
\end{figure}

Equation~\ref{eq:gravity} represents the gravitational model.

\begin{equation}
f(T)=T+\alpha T^{2}
\label{eq:gravity}
\end{equation}

\section{References}

For more details, please refer to \cite{solar1} and \cite{solar2}.

\bibliographystyle{IEEEtran}
\bibliography{ref}

\end{document}